% =====================================================================
% Build de preview -- Sectiunile 1 (Introduction) si 2 (Related Work).
% Nu e fisierul de submisie: nu contine abstract, Approach, Results,
% Limitations. Corpul e paste_intro.tex integral.
%
% De la revizia 6 include si Related Work, ca sa fie vizibile subsectiunile
% 2.1-2.5. Pana atunci se taia inainte de \section{Related Work}, si de aceea
% sec:related era printre etichetele fixate manual mai jos; acum e o eticheta
% reala si a fost scoasa de acolo (altfel: "Label multiply defined").
% =====================================================================

\documentclass[11pt]{article}

\usepackage[review]{acl}

\usepackage{times}
\usepackage{latexsym}
\usepackage[T1]{fontenc}
\usepackage[utf8]{inputenc}
\usepackage{microtype}
\usepackage{inconsolata}
\usepackage{graphicx}

\usepackage{booktabs}
\graphicspath{{figures/}}

\title{[titlul se decide ulterior]}

\author{}

\begin{document}
\maketitle

% Abstractul e inclus DOAR ca paginarea sa fie realista. Fara el, pagina 1 are
% o coloana in plus libera si masuratoarea "incape in 2 pagini" ar fi falsa.
% =====================================================================
% PASTE FILE -- abstract nou, scris pentru introducerea de la revizia 2.
% Se copiaza peste blocul \begin{abstract}...\end{abstract} din
% acl_latex.tex in Overleaf.
%
% Cel vechi (report/report.tex) era de raport de curs: descria ce s-a
% facut, fara niciun rezultat. Asta il inlocuieste. ~140 de cuvinte.
% Anonimizat: fara nume, fara link.
% =====================================================================

\begin{abstract}
Lexical complexity prediction is normally solved by fine-tuning an
encoder on annotated data in the target language, which excludes every
language that has none. We predict the lexical complexity of Romanian
words in context without using a single Romanian training label, by
fusing four signals that carry complexity knowledge from outside the
language: cross-lingual transfer through a frozen multilingual encoder
trained on English and Romance annotations, a complexity rating and
token surprisal from a Romanian language model, and a reading-effort
score derived from eye-tracking. The system reaches $69.74$ on the
shared-task metric, against $10.27$ for surface features alone, and
Pearson $r = 0.913$ on held-out Romanian items. The largest single gain
($+4.09$) comes from a lemma-frequency correction that repairs the
frequency estimates inflection corrupts, and transfers to any
morphologically rich language served by the same frequency lists. We
also report two negative results: with no target labels, fine-tuning
the encoder is consistently worse than freezing it, and in-source
validation does not predict transfer quality.
\end{abstract}


% Referintele inainte catre sectiuni care nu exista in acest preview sunt
% fixate manual la numerele pe care le au in paperul complet, ca sa nu apara
% "Section ??". Numai pentru preview -- nu se lipeste in Overleaf.
\makeatletter
\def\pinlabel#1#2{\expandafter\gdef\csname r@#1\endcsname{{#2}{1}{}{}{}}}
\pinlabel{section:approach}{3}
\pinlabel{sec:tuning}{3.5}
\pinlabel{sec:results}{4}
\pinlabel{sec:14items}{4.2}
\pinlabel{sec:negative}{5}
\pinlabel{section:conclusions}{6}
\makeatother

% =====================================================================
% PASTE FILE -- introducere + related work, scrise pentru MRL 2026.
% Se copiaza peste \section{Introduction} din acl_latex.tex in Overleaf.
% Intrarile .bib necesare sunt in paste_bib_additions.bib (toate verificate
% pe ACL Anthology).
%
% Structura urmeaza modelul hgs (cs.columbia.edu/~hgs/etc/intro-style.html):
%   P1 motivatie -> P2 problema specifica -> P3 contributii ->
%   P4 diferentiere -> P5 roadmap
% Versiune anonimizata (MRL e double-blind): fara link de GitHub, fara nume.
%
% Revizia 2: fiecare afirmatie are acum ori o citare, ori un numar din paper.
% Numerele noi din P2 (11.7 / 36.0 / 0.86) sunt calculate pe setul de test,
% n=1587 -- vezi sectiunea Approach.
%
% Revizia 3 (27 iulie 2026): pornita de la revizia 2, remodelata retoric dupa
% paperul-exemplu 2605.07466 (Anghel et al., fatty pancreas), pastrand
% structura hgs. Ce s-a schimbat:
%   - motivatia despicata in doua paragrafe: miza (P1) -> reteta standard si
%     presupozitia ei (P2). Problema specifica devine P3, contributiile P4,
%     diferentierea P5, roadmap-ul P6.
%   - contributiile: enumerate cu itemi de 5-8 randuri -> itemize cu lead-in
%     bold si 3-5 randuri, o revendicare si o cifra fiecare.
%   - paragraful separat despre tuning (rev. 2) comprimat la o fraza dupa
%     lista de contributii; tratarea integrala ramane in sec:tuning si in
%     Limitations, unde era deja.
%   - P5 diferentiere scurtata; citarile fine coborate in Related Work.
%
% ATENTIE, verificat pe ACL Anthology pe 27 iulie 2026: shared task-ul CWI
% 2018 (yimam-etal-2018-report) a avut patru piste -- EN/DE/ES monolingv PLUS
% "a multilingual track with a French test set", adica exact un setup zero-shot
% cross-lingual. Deci NU se poate revendica "primul sistem LCP fara etichete in
% limba tinta". Claim-ul de noutate e limitat la romana, iar pista franceza e
% acum recunoscuta explicit in Related Work si diferentiata (clasificare binara
% si probabilistica, nu predictie continua; date de antrenare in trei limbi,
% una romanica).
%
% Intrari .bib noi fata de revizia 2: qiu-etal-2024-complex,
% saggion-etal-2024-lexical. Total referinte: 18.
%
% Revizia 4 (27 iulie 2026): reformulari cerute la review, in P1, P2 si P6.
% Nicio schimbare de continut, de cifre sau de citari.
%   - P1, prima fraza: "Text is only as accessible as its hardest words"
%     (aforism, neclar) -> deschidere pe cine e afectat, "Children, language
%     learners, and readers with limited literacy or a reading disability are
%     held up by the same thing: individual words they do not know." Aleasa
%     dintre trei variante propuse: mizele umane (asta), cele trei aplicatii
%     comprimate intr-o singura intrebare, si o fraza scurta fara figuri de stil.
%     O varianta intermediara ("Whether a text can be read by the audience it
%     was written for depends on which of its words that audience finds
%     difficult") a fost si ea respinsa, ca prea incarcata de subordonate.
%   - P1: scos "Automatic simplification" -> "text simplification"; in NLP
%     simplificarea e oricum automata, adjectivul nu adauga nimic.
%   - P2: scos "The field has converged on a single recipe" -> afirmatia
%     directa "Nearly every LCP system is built the same way: ...".
%   - P2: scos "A decade of shared tasks" -> "the shared tasks in the area
%     have established", fara cuantificarea in ani.
%   - P6 (roadmap): rescris. Ghidul cere explicit "Avoid redundant phrasing,
%     'In Section 2, In section 3, ... In Section 4, ...'", iar versiunea
%     anterioara facea exact asta pe patru sectiuni la rand. Acum enumera
%     continutul, nu sectiunile, si nu mai anunta Conclusions.
%
% Verificat la revizia 4, fara modificari necesare: concordanta timpurilor in
% P2 (singurul perfect compus, "have established", e corect - proces incheiat
% cu relevanta in prezent) si fraza-tema a lui P3, care leaga de P2 prin "such
% a language" si anunta cele doua probleme pe care paragraful le livreaza apoi.
%
% Revizia 5 (27 iulie 2026), tot la cererea lui Alexandru. Deschiderea a fost
% remodelata dupa introducerea din 2605.07466_6048.pdf, care e paperul-model
% agreat, pastrand structura din ghid.
%   - P1 rescris integral. Modelul isi deschide introducerea numind si definind
%     obiectul de studiu, cu citare, apoi trece la consecinte concrete; nicio
%     figura de stil. Deci P1 incepe acum cu definitia LCP si continua cu cele
%     trei aplicatii care depind de ea. Respinse pe rand, ca "vrajeala": varianta
%     cu aforismul ("Text is only as accessible as its hardest words"), cea cu
%     subordonate ("Whether a text can be read by the audience...") si cea cu
%     mizele umane ("Children, language learners, and readers with limited
%     literacy..."). A cazut si formularea "The systems built to help them".
%   - SCOS TOT em-dash-ul din introducere, sase locuri: P1, P3, doua in lista de
%     contributii si unul in paragraful de tuning. Inlocuit cu virgule, punct si
%     virgula sau reformulare. Modelul nu foloseste em-dash deloc. De pastrat
%     regula asta la orice text nou.
%   - P3: scoasa fraza "Fourteen labelled Romanian target words exist, enough to
%     sanity-check a system but neither to train one nor to validate one...".
%     Faptul portant (nu exista set de dezvoltare) a ramas, dar reformulat, iar
%     paragraful anunta acum explicit doua obstacole si le livreaza in ordine.
%     Cele 14 exemple sunt introduse acolo unde chiar sunt folosite, in prima
%     contributie, formulate ca sa se sustina singure.
%
% Structura ramane 1. Introduction + 2. Related Work: ghidul o cere explicit la
% P4 ("you can refer to a future section where specific details and differences
% will be given"), iar paperul-model face la fel (1 INTRODUCTION, 2 RELATED WORK
% cu subsectiuni). Verificat pe 27 iulie 2026.
%
% Revizia 6 (27 iulie 2026): aliniere completa la paperul-model, cerute amandoua.
%   - Lista de contributii trecuta pe verb la persoana I, ca in model ("We build",
%     "We introduce", "We show", "We report"), si scos lead-in-ul ingrosat.
%     Modelul nu foloseste bold in lista de contributii. Continutul si cifrele
%     sunt neschimbate, doar constructia frazei.
%   - Cele cinci \paragraph{} din Related Work promovate la \subsection{}, cu
%     titluri in title case: 2.1 Lexical Complexity Prediction, 2.2 Cross-Lingual
%     Transfer Without Target Labels, 2.3 Surprisal and Eye Movements as
%     Difficulty Signals, 2.4 Evaluation Without a Development Set, 2.5 Romanian
%     Resources. Modelul are Related Work cu subsectiuni numerotate (2.1, 2.2).
%   - Scos em-dash-ul si din Related Work, inca cinci locuri (2.1, 2.2, 2.3 si
%     doua in 2.4). La revizia 5 fusese curatata doar introducerea, pentru ca
%     Related Work nu intra atunci in compilarea de preview. Acum fisierul e
%     curat integral.
%
% Revizia 7 (27 iulie 2026): consecinte ale numerotarii subsectiunilor.
% Masurate, subsectiunile ieseau 155 / 125 / 86 / 60 / 35 de cuvinte; ca
% \paragraph{} dezechilibrul nu se vedea, ca subsectiuni numerotate, da.
%   - 2.5 Romanian Resources era un ciot de 35 de cuvinte care amesteca RoLlama3
%     (resursa romaneasca reala) cu SHAP (metoda de interpretabilitate, fara
%     legatura cu titlul). SHAP a fost scos de aici; subsectiunea are acum o teza
%     proprie (romana e bine servita de resurse generale si deloc de resurse
%     specifice sarcinii, iar designul nostru decurge din asta) si ~120 de cuvinte.
%   - ATENTIE: scoaterea lui SHAP de aici lasa \citep{lundberg2017shap} orfan,
%     pentru ca paste_body.tex foloseste SHAP de sapte ori fara sa-l citeze
%     niciodata. Citarea a fost mutata in paste_body.tex, la prima folosire
%     ("A SHAP analysis \citep{lundberg2017shap} over named features"). Fara
%     mutarea asta, referinta ar fi disparut din bibliografia lucrarii complete.
%   - 2.1 deschidea cronologic ("CWI was originally posed as..."), nu cu teza
%     subsectiunii. Acum incepe cu "Annotated complexity data exists for only a
%     handful of languages, and Romanian is not among them", iar istoricul vine
%     dupa, ca sprijin. Cerinta de fraza-tema din ghid.
%   - Acronimul CWI era folosit in paragraful de diferentiere si in 2.1 fara sa
%     fie definit nicaieri. Expandat acum la prima aparitie, in P2.
%
% =====================================================================
% Revizia 8 (27 iulie 2026). DOUA lucruri diferite, de nu confundat.
%
% (A) TREI AFIRMATII FALSE, SCOASE. Nu erau chestiuni de stil.
% Setul nostru de test de 1587 de tinte este RoLCP, publicat in RALS (Anghel,
% Cristea, Creanga, Nisioi), EMNLP 2025 main. Sergiu Nisioi, ultimul autor, este
% coordonatorul lui Alexandru. Descoperit plecand de la paperul-model ETRA
% (Kuvshynova et al. 2026), care citeaza RoLCP. Dimensiunea se potriveste exact
% si formularea e tot continua in [0,1].
%   - "To our knowledge, no prior system predicts lexical complexity for
%     Romanian" (contributia 1): STERSA. Ei raporteaza un Ridge supervizat cu
%     r = 0.71 pe exact aceste date.
%   - "Romanian has no complexity-annotated training set and no development set"
%     (P3): REFORMULATA. Datele exista (3921 de exemple: 569 HT + 1765 WT +
%     1587 RoLCP, public, CC BY-NC-SA). Noi ALEGEM sa nu le folosim, ceea ce e
%     intelegerea cu coordonatorul. P3 declara acum constrangerea ca alegere de
%     design, nu ca lipsa a datelor din lume.
%   - "we are aware of no published LCP system for it" (2.1): STEARSA.
%   - "the fourteen labelled Romanian items that exist" -> "released with the
%     shared task". Cele 14 sunt ce s-a dat participantilor.
%   Aceleasi corectii facute si in paste_body.tex (Conclusions, "the only
%   Romanian labels available") si paste_limitations.tex ("the entirety of the
%   Romanian supervision that exists").
%
%   Ce NU s-a schimbat, pentru ca e adevarat: nu folosim nicio eticheta de
%   antrenare romaneasca. Asta ramane teza lucrarii, in abstract si peste tot.
%
%   % TODO, de decis cu Sergiu inainte de submisie:
%     1. Se citeaza RALS? Alexandru a cerut explicit sa NU fie citat deocamdata,
%        deci nu e. Dar lucrarea nu citeaza in acest moment NICIO sursa pentru
%        setul de 1587 de exemple, iar ACL cere proveninta datelor.
%     2. Se citeaza Kuvshynova et al. 2026 (ETRA)? Combina RoLCP cu eye-tracking
%        MultiplEYE romanesc, adica exact unghiul nostru. E cea mai apropiata
%        lucrare de a noastra si e necitata.
%     3. Daca (1) e da: se citeaza si r = 0.71 al lor ca punct de comparatie?
%        Atentie, protocoalele difera (ei cross-valideaza supervizat, noi
%        prezicem fara antrenare), deci trebuie spus explicit.
%
% (B) REFORMULARILE CERUTE de Alexandru la review.
%   - P1 rescris a treia oara. "Lexical complexity prediction (LCP) is the task
%     of estimating..." suna a definitie de dictionar, "se vede ca nu e scrisa de
%     om". Acum urmeaza MISCAREA din introducerea ETRA (de ce conteaza -> cine e
%     afectat -> numele sarcinii -> ce permite), fara sa semene 1 la 1 cu ea.
%   - "but Romanian is heavily inflected and inflected forms are systematically
%     missing from frequency lists" -> "Romanian inflects heavily, and the
%     frequency lists available for it cover inflected forms poorly".
%   - Scos "by 0.86 Zipf units on average" din intro. Cele doua procente raman,
%     pentru ca justifica cele doua ramuri ale corectiei pe lema, dar detaliul cu
%     unitatile Zipf era duplicat verbatim din Approach.
%   - "A method for Romanian has to transfer complexity knowledge from outside
%     the language and repair the frequency signal it transfers on" era neclar.
%     Spart in doua cerinte explicite: "bring in complexity knowledge from
%     languages that have it, and repair the frequency estimates before leaning
%     on them".
%   - Scos acronimul din "We build an LCP system"; contributia 1 incepe acum cu
%     "We combine four label-free signals into a single system".
%
% RAMAS DE FACUT, semnalat: paste_body.tex si paste_limitations.tex inca folosesc
% em-dash. Regula "fara em-dash" a fost aplicata doar pe paste_intro.tex, pentru
% ca livrabilul curent e introducerea.
% =====================================================================

\section{Introduction}
\label{section:intro}

Making a text usable by a given audience depends, before anything else,
on its vocabulary. A word that is routine for a specialist can stop a
language learner or a reader from outside the field, taking the
surrounding sentence with it. Deciding which
words will cause this, and for whom, is what lexical complexity
prediction (LCP) does: it assigns a word in context a score on a
continuous scale \citep{shardlow-etal-2021-semeval}. Text simplification
uses those scores to choose what to replace, readability assessment to
aggregate difficulty over a document, and language-teaching material to
keep vocabulary within a learner's range; complexity prediction is for
this reason typically the first stage of a simplification pipeline
\citep{qiu-etal-2024-complex}.

Nearly every LCP system is built the same way: a pretrained encoder is
fine-tuned on complexity annotations in the language of interest. The
recipe works wherever those annotations exist, and it quietly assumes the
one resource most languages do not have. Surveying the datasets published
between 2018 and 2024, \citet{shardlow-etal-2024-multilingual} find nine
of them, covering eight languages. A language outside that set has
nothing to fine-tune on, and buying its way in means paying for a new
annotation effort of its own.

Our target language is Romanian, and we adopt the hardest version of the
setting: no Romanian complexity label enters the system at any point, for
training, validation or model selection. That constraint creates two
obstacles. The first is that no design decision
can be checked against the target language, so every choice has to be
justified from elsewhere. The second is that Romanian degrades the single
most reliable predictor of processing difficulty. Word frequency governs how long readers fixate
a word \citep{rayner1998eye} and it is the top-ranked feature in our own
attribution analysis (Section~\ref{sec:results}), but Romanian inflects
heavily, and the frequency lists available for it cover inflected forms
poorly: of the 1587 Romanian target tokens in our test set, $11.7\%$ have
no entry at all and would be scored as maximally rare, and for a further
$36.0\%$ the listed form is rarer than its own lemma. Ordinary words are
therefore made to look rare, and the words this happens to are precisely
the morphologically marked ones. A method that works here has to do two
things at once: bring in complexity knowledge from languages that have
it, and repair the frequency estimates before leaning on them.

In this paper, we predict the lexical complexity of Romanian words in
context without a single Romanian training label, and we make four
contributions. First, we combine four label-free signals into one
system: cross-lingual transfer through a frozen multilingual encoder, a
language-model complexity rating, surprisal, and a reading-effort score
derived from eye-tracking. It reaches $69.74$ on the shared-task metric,
against $41.76$ when restricted to the signals the shared task itself
provides and $10.27$ for surface features alone, and Pearson
$r = 0.913$ on the fourteen labelled Romanian items released with the
task. Second, we introduce a lemma-frequency correction for
morphologically rich languages, substituting the lemma frequency when the
inflected form is missing or rarer; it is worth $+4.09$, the largest
single gain in the system, and applies to any inflected language served
by the same frequency lists. Third, we show why the four signals compose
rather than duplicate each other: the cross-lingual signal is essentially
a frequency detector, the rating and surprisal are near-independent
($r = 0.07$), and eye-tracking contributes through its \emph{late}
measures the contextual-integration difficulty the others miss.
Fourth, we report two negative results: fine-tuning the encoder, full or
LoRA, is consistently \emph{worse} than leaving it frozen ($0.461$ and
$0.462$ against $0.515$ in-source), and in-source validation is an
unreliable guide, since three configurations that win on
English-to-Romance transfer lose on Romanian.

Our setting differs from the LCP literature in the resource it assumes
rather than in the model it uses: where that literature trains on
target-language annotations, we treat complexity as knowledge to be
imported. Two smaller inversions follow: we read gaze as evidence about
complexity rather than predicting gaze from text, and we find the frozen
encoder to beat fine-tuning, the opposite of the usual cross-lingual
recipe. Section~\ref{sec:related} sets all three against the specific
work they concern.

The rest of the paper describes the four signals and their fusion
(Section~\ref{section:approach}), reports the ablation and the analysis
(Section~\ref{sec:results}), and collects what did not work
(Section~\ref{sec:negative}).

% =====================================================================

\section{Related Work}
\label{sec:related}

\paragraph{Lexical complexity prediction.}
Annotated complexity data exists for only a handful of languages, and
Romanian is not among them. The task itself began as classification over
annotated target words \citep{yimam-etal-2018-report}; SemEval-2021
Task~1 replaced it with a continuous scale on the CompLex corpus
\citep{shardlow-etal-2021-semeval}, which is the formulation we adopt.
Zero-shot evaluation is not itself new: the 2018 task already had a
multilingual track whose French test set came with no French training
data, though it asked only for binary classification and supplied
training data in three languages, one of them Romance. Coverage has since
grown to
Catalan, Spanish, Mandarin and Cantonese
\citep{saggion-etal-2024-lexical, qiu-etal-2024-complex} and to the ten
languages of the MLSP shared task \citep{shardlow-etal-2024-bea}, but the
review of \citet{shardlow-etal-2024-multilingual} shows annotated LCP
data still concentrated in a few high-resource languages. Romanian
appears in none of these resources. We use the English and Romance
portions of them as the source side of our transfer.

\paragraph{Cross-lingual transfer without target labels.}
Multilingual encoders such as XLM-R \citep{conneau2020xlmr} place
translation-equivalent words close together, which makes zero-shot
transfer of a regressor trained on one language plausible. Transfer
quality is far from uniform across languages, however, and degrades for
languages with less pretraining data and richer morphology
\citep{pires-etal-2019-multilingual, wu-dredze-2020-languages}, which is
precisely the regime Romanian occupies. Transfer is moreover normally
realised by fine-tuning the encoder on the source language; that feature
extraction can outperform fine-tuning when the target task is distant
from the source is known \citep{peters-etal-2019-tune}, but we are aware
of no prior report of it for cross-lingual complexity transfer. This
motivates our frozen-encoder design, and the lemma-frequency correction
completes it: the encoder supplies cross-lingual semantics, the
correction repairs the lexical feature that morphology corrupts.

\paragraph{Surprisal and eye movements as difficulty signals.}
Surprisal theory predicts that processing effort scales with the
negative log probability of a word in context \citep{hale2001probabilistic},
and gaze measures such as longer fixations, re-reading and regressions
are the standard behavioural correlate of that effort
\citep{rayner1998eye}. Language-model surprisal predicts reading times
across languages \citep{hollenstein-etal-2021-multilingual}, including
in Romance languages \citep{alves-2025-benchmarking}. Both literatures
run from text to behaviour; we take the converse view, reading surprisal
and gaze as label-free evidence about complexity.

\paragraph{Evaluation without a development set.}
Selecting a model against the same data used to report its score inflates
that score \citep{gorman-bedrick-2019-need}. With no development set to
tune on, we state our selection procedure explicitly and quantify its
effect (Section~\ref{sec:tuning}).


\bibliography{custom}

\end{document}
